%% --------------------------------------------------------
\section{Coupled Cluster методи (CCSD, CCSD(T))}
%% --------------------------------------------------------

%% --------------------------------------------------------
\subsection{Теоретичні основи Coupled Cluster}
%% --------------------------------------------------------

Coupled Cluster (CC) --- один з найточніших методів для врахування електронної кореляції. На відміну від CI, CC є size-extensive (правильно масштабується для великих систем).

Хвильова функція у CC записується через експоненціальний оператор збудження:

\begin{equation}
    |\Psi_{CC}\rangle = e^{\hat{T}} |\Phi_0\rangle
\end{equation}
%
де $|\Phi_0\rangle$ --- HF детермінант, а $\hat{T}$ --- оператор кластерів:

\begin{equation}
    \hat{T} = \hat{T}_1 + \hat{T}_2 + \hat{T}_3 + \ldots
\end{equation}

\paragraph{CCSD (Coupled Cluster Singles and Doubles)}
Враховує тільки $\hat{T}_1$ (одиночні збудження) та $\hat{T}_2$ (подвійні збудження):

\begin{align}
    \hat{T}_1 & = \sum_{i}^{\text{occ}} \sum_{a}^{\text{virt}} t_i^a \, \hat{a}^\dagger_a \hat{a}_i                                                 \\
    \hat{T}_2 & = \frac{1}{4} \sum_{ij}^{\text{occ}} \sum_{ab}^{\text{virt}} t_{ij}^{ab} \, \hat{a}^\dagger_a \hat{a}^\dagger_b \hat{a}_j \hat{a}_i
\end{align}

\paragraph{CCSD(T) --- <<золотий стандарт>>}
Додає пертурбативну корекцію для потрійних збуджень $\hat{T}_3$:

\begin{equation}
    E_{CCSD(T)} = E_{CCSD} + E_{(T)}
\end{equation}

де $E_{(T)}$ обчислюється за формулами четвертого порядку теорії збурень.


%% --------------------------------------------------------
\subsection{CCSD розрахунок атома Гелію}
%% --------------------------------------------------------

\inputcode{code_8.py}

%% --------------------------------------------------------
\subsection{CCSD для відкритих оболонок (UCCSD)}
%% --------------------------------------------------------

\inputcode{code_9.py}

%% --------------------------------------------------------
\subsection{Порівняння ієрархії методів}
%% --------------------------------------------------------

\inputcode{code_10.py}

%% --------------------------------------------------------
\subsection{T1 діагностика та багатоконфігураційний характер}
%% --------------------------------------------------------

T1 діагностика --- індикатор того, наскільки система є одноконфігураційною:

\begin{equation}
    T_1 = \frac{\|t_1\|}{\sqrt{N_{\text{elec}}}}
\end{equation}

\inputcode{code_11.py}

%% --------------------------------------------------------
\subsection{Обчислювальна складність та масштабування}
%% --------------------------------------------------------

\begin{table}[h]
    \centering
    \caption{Порівняння обчислювальної складності методів}
    \label{tab:cc_complexity}
    \begin{tabular}{lccc}
        \hline
        \textbf{Метод} & \textbf{Масштабування} & \textbf{Пам'ять}   & \textbf{Disk}      \\
        \hline
        HF             & $\mathcal{O}(N^4)$     & $\mathcal{O}(N^2)$ & Мало               \\
        MP2            & $\mathcal{O}(N^5)$     & $\mathcal{O}(N^4)$ & $\mathcal{O}(N^4)$ \\
        CCSD           & $\mathcal{O}(N^6)$     & $\mathcal{O}(N^4)$ & $\mathcal{O}(N^4)$ \\
        CCSD(T)        & $\mathcal{O}(N^7)$     & $\mathcal{O}(N^4)$ & $\mathcal{O}(N^4)$ \\
        CCSDT          & $\mathcal{O}(N^8)$     & $\mathcal{O}(N^6)$ & $\mathcal{O}(N^6)$ \\
        FCI            & $\mathcal{O}(e^N)$     & $\mathcal{O}(e^N)$ & $\mathcal{O}(e^N)$ \\
        \hline
    \end{tabular}
\end{table}

де $N$ --- характерний розмір системи (базисних функцій або електронів).
